% Include-only supplement. Requires amsmath and amssymb.
% The finite source checker is audit_command_permutation.py.
\begingroup
\newcommand{\Kp}{\mathsf P}
\newcommand{\Ki}{\mathsf I}
\newcommand{\Kq}{\mathsf Q}
\newcommand{\Kd}{\mathcal D}
\newcommand{\Kax}{\mathsf{AX}}
\newcommand{\Kmp}{\mathsf{MP}}
\newcommand{\Kdrop}{\mathsf{DROP}}
\newcommand{\Knip}{\mathsf{NIP}}
\newcommand{\Kenc}{\operatorname{enc}}
\newcommand{\Kblk}{\mathcal B}
\newcommand{\Kdiscard}{\mathsf{Discard}}

\section{All-natural calibration of the permuted command source}
\label{sec:k216-source-calibration}

The source-level change that makes the new initialization possible is a
permutation of the three small command selectors. The arithmetic
commands and the dominated axiom family are unchanged. This section
proves that the permutation preserves the full consistency calibration,
including raw candidate numbers, the zero alias, arbitrary formula
codes, and the actual source seed with next-candidate register $4$.
It does not infer a binary runtime theorem from a source subtree count;
the binary runtime and blank-tape certificates are separate obligations.

\subsection{Exact finite source identity}

Write
\[
 \Kp(a,b)=\frac{(a+b)(a+b+1)}2+b,\qquad
 \Ki(a,b)=4\Kp(a,b),\qquad
 \Kq_i a=4\Kp(i,a)+2.
\]
The native truth and falsity codes are respectively $0$ and $1$.
In particular, $\Ki(0,0)=0$. Let $\Kd$ be the dominated numerical
calculus of the preceding dominated-rho8 construction: truth, numerical
modus ponens, theorem-only generalization, the seven retained logical
rows, the gated theory bundle, and
\[
 L(A,B,C,U)=
 \Ki\bigl(\Ki(\Ki(C,U),A),
            \Ki(\Ki(A,C),\Ki(B,C))\bigr),\qquad U\leq A.
\]
The inherited finite propositional certificate proves equality of the
$\Kd$ and K219 theorem sets for \emph{every natural code}. Together with
the inherited numerical-calculus-to-ZFC calibration, it gives
\begin{equation}
 \Kd\vdash1\quad\Longleftrightarrow\quad
 \exists p\;\operatorname{Proof}_{\mathrm{ZFC}}(p,\bot).
 \label{eq:k216-inherited-calibration}
\end{equation}
No well-formedness or successful formula-decoding test is inserted into
this equivalence.

The source checker first reconstructs the 205-subtree noncommon-input
reference directly from the pinned K219 source. Its independently
implemented structured-control and address-counter interpreters agree
on every primitive and every branch outcome. The resulting guarded
ordinal control graph is exactly that of the frozen 202-subtree source
with SHA-256
\begin{center}\ttfamily\footnotesize
abb7229877759170b61dfcfe21ca0889b3df2f91f4215383817a125eb0f8996c.
\end{center}
Thus the frozen source has the arithmetic command semantics proved in
the dominated-rho8 report, with counter/cache field order YX.

In this frozen source, the MP, DROP, and NIP bodies occur at paths
$111$, $110101$, and $110100101$, respectively. They contain
$106$, $26$, and $26$ primitive occurrences. Each has zero escaping
break and continuation demand. NIP contains the halting primitive;
all its nonhalting exits are normal returns. The branch-skipping
instructions outside the three bodies remain fixed when the bodies
are permuted. The checker traces their ordinary continuations: the
last body returns directly to source ordinal $0$, the middle body
returns through the unconditional break at $4724$, and the innermost
body returns through the breaks at $4697$ and $4724$. Each route ends
at ordinal $0$ without another arithmetic command.

For the source denoted $210$, put NIP, DROP, and MP in these three
selector positions. For $120$, put DROP, NIP, and MP there instead.
The exact selector partitions and wire formats are
\[
\begin{array}{c|c|c|c|c}
 &\text{selector }0&\text{selector }1&\text{selector }2
 &\text{selector }c+3\\\hline
210&\Knip&\Kdrop&\Kmp&\Kax(c;x,y)\\
120&\Kdrop&\Knip&\Kmp&\Kax(c;x,y)
\end{array}
\]
and
\begin{align*}
 f_{210}(\Knip)&=[0],& f_{210}(\Kdrop)&=[1],&
 f_{210}(\Kmp(x))&=[2,x],\\
 f_{120}(\Knip)&=[1],& f_{120}(\Kdrop)&=[0],&
 f_{120}(\Kmp(x))&=[2,x],\\
 f_\pi(\Kax(c;x,y))&=[c+3,y,x]&& (\pi=210,120).
\end{align*}
The semicolon in $\Kax(c;x,y)$ separates the selector parameter from
the counter increment $x$ and cache increment $y$; it does not denote
wire order.

The independent checker proves these partitions directly from the three
selector decrement tests, at primitive ordinals $23,24,25$. In the
$210$ source, their zero branches lead to NIP at ordinal $4725$, DROP
at $4698$, and MP at $4591$. Three positive outcomes enter AX at ordinal
$26$, with residual selector $c$. This last case is the symbolic case
of every input selector at least $3$, not a finite sampling of large
selectors.

Both sources have 202 distinct subtrees, 4,751 primitive occurrences,
6,176 guarded primitive cases, maximum address depth 44, and maximum
zero run 6. All 4,593 primitive occurrences exterior to the replaced
bodies, including the entire AX subtree, remain literal. The checker
also verifies that each candidate tree is \emph{exactly} the indicated
permutation of the frozen tree. The candidate source hashes are
\begin{center}\footnotesize\ttfamily
210: 8b5dc46f10d927b2db5e4bcb522d482779ec61c72c82717d7bf12fca3550a3bc\\
120: ee1f7c7cb216a80315449371b040ee094728a655439954225ed4be4635296680.
\end{center}

\subsection{Command equations and soundness for raw codes}

At a command boundary write the retained tuple as $(t,w,d,h)$: top,
encoded stack, additive counter, and cached formula. The notation
$[x_1,\ldots,x_k;w]$ means right-nested pairing, and
$\operatorname{unpair}(0)=(0,0)$. The exact command equations, independent
of which small tag names the command, are
\begin{align*}
 \Kmp(x)(t,w,d,h)&=
 \begin{cases}
 (d+x,\Kp(0,w),d+x,0),&\Ki(t,d+x)=h,\\
 (0,w,d+x,0),&\Ki(t,d+x)\ne h;
 \end{cases}\\
 \Kdrop(t,w,d,h)&=(h,\Kp(t,w),d,h),\\
 \Knip(t,w,d,h)&=
 \begin{cases}
 \mathsf{halt},&h=1,\\
 (t,w',d,h'),&(h',w')=\operatorname{unpair}(w),\ h\ne1.
 \end{cases}
\end{align*}
For $\Kax(c;x,y)$ put
\[
 U=d,\qquad A=d+x,\qquad B=h+y,\qquad C=c,
\]
and let $G(A,B,c)$ be zero when $c=A$ or $c=B$, and the retained theory
bundle otherwise. In the inherited stack order its output is
\begin{align*}
 t'&=G(A,B,c),& d'&=0,&h'&=0,\\
 w'&=[R_1(A,B,c),\ldots,R_7(A,B,c),L(A,B,c,U),\Kq_Uh,t;w].
\end{align*}
Here the $h$ in $\Kq_Uh$ is the incoming cache, before adding $y$.
The bound $U\leq A$ holds for every raw payload because $A=d+x$.

A body with no escaping break or continuation has the same effect at
any of the three command positions. The branch wrappers route every
normal return back to the command loop. Consequently the equations
above hold for the permuted sources on every natural input tuple,
including payloads arising from incomplete or otherwise unintended
candidate encodings.

\paragraph{Soundness invariant.}
If $t$, $h$, and every successive head of $w$ are $\Kd$-theorems,
every nonhalting command preserves that property. A successful MP is
exact numerical modus ponens; a failed MP inserts truth $0$. DROP and
NIP only move theorem codes. AX inserts retained axioms, a valid
dominated-row instance, a generalization of a theorem, and either a
permitted bundle or truth. The only halt is the NIP test $h=1$.
Therefore halting from any such safe seed implies $\Kd\vdash1$.
This argument quantifies over all natural codes and does not discard
``junk'' candidates.

\subsection{An explicit compiler retained under the tag change}

For completeness, a strong block $\Kblk_X$ must, from every $(t,w,0,0)$,
either halt or return $(X,w,0,0)$. The command lists below use command
\emph{names}; their concrete wire encodings use $f_{210}$ or $f_{120}$.
Concatenation denotes sequential execution, and a superscript on a
command denotes repetition.

Define
\[
 \Kdiscard=[\Kmp(1),\Kax(1;0,0)]\Knip^{10},\qquad
 J(u)=\begin{cases}[],&u=0,\\{}[\Kmp(u)],&u>0.\end{cases}
\]
The first MP fails because $\Ki(t,1)>0$, whereas the incoming cache is
zero. The following AX has $U=A=c=1$ and a zero gate. Its nine newly
generated codes are all different from $1$, and its saved old top is
zero. Hence ten NIPs return exactly $(0,w,0,0)$.
The block $J(u)$ maps $(0,w,0,0)$ to $(0,w,u,0)$: if $u>0$,
$\Ki(0,u)>0$ forces failure. Omitting MP when $u=0$ avoids the extra
stack zero that a successful $\Kmp(0)$ would insert.

Truth uses $\Kdiscard$. A dominated-row leaf with $U\leq A$ uses
\[
 \Kdiscard\;J(U)\;[\Kax(c;A-U,B)]\;\Knip^8\;[\Kdrop]\;\Knip^3.
\]
For retained row $R_j(A,B,c)$, where $1\leq j\leq7$, use
\[
 \Kdiscard\;[\Kax(c;A,B)]\;\Knip^j\;[\Kdrop]\;\Knip^{11-j}.
\]
For a permitted bundle use
$\Kdiscard[\Kax(c;A,B)]\Knip^{10}$.
These suffixes select the indicated top or stack entry and restore the
underlying stack and zero cache. After expanding the negation and
existential abbreviations, all eight row codes are implications,
and $\Kq_Uh\equiv2\pmod4$, so none of these generated stack codes is
falsity $1$.

Given a block for $X$, a block for $\Kq_iX$ is
\[
 \Kblk_X\;[\Kdrop]\;J(i)\;[\Knip,\Kax(i;0,0)]\;
 \Knip^9\;[\Kdrop,\Knip,\Knip].
\]
After DROP and $J(i)$, the current top and cache are zero, the counter
is $i$, and $X$ is the first stack head. NIP places $X$ in the cache.
AX then has $U=A=c=i$, so its top gate is zero and its ninth generated
stack entry is exactly $\Kq_iX$. The selection suffix returns that code
with the original stack and zero counter and cache.

Given blocks for $X$ and $\Ki(X,Y)$, use
\begin{multline*}
 \Kblk_X\;[\Kdrop]\;\Kblk_{\Ki(X,Y)}\;
 [\Knip,\Kdrop,\Knip,\Kmp(Y),\Kax(Y;0,0)]\\
 {}\Knip^{10}\;[\Kdrop,\Knip,\Knip].
\end{multline*}
Unless an earlier NIP has halted, the two premises reach the top and
cache in the required order, and MP succeeds. The reset AX has
$U=A=c=Y$ and a zero gate. The final selection returns $Y$ with the
original stack and zero counter and cache, or halts if a tested cache
is already $1$.

Induction on a finite $\Kd$ derivation proves the strong block contract
for every derived natural code. The construction is primitive recursive
on the derivation. The new selectors alter only the wire names of these
commands, so every calculation above applies verbatim to both permuted
sources.

\subsection{Exact packing and cofinality from the actual cutoff}

For a finite field list $fs$ and tail $r$, write $E(fs,r)$ for its
right-folded Cantor pairing. Command tails are encoded by
\[
 \Kenc_\pi([])=0,\qquad
 \Kenc_\pi(k::L)=1+E(f_\pi(k),\Kenc_\pi(L)).
\]
A newly enumerated candidate is loaded directly from its number. A
subsequent command is decoded from one less than the current positive
proof-list value. Thus the number for a whole first candidate
$k::L$ is $E(f_\pi(k),\Kenc_\pi(L))$, without an outer successor.
This distinction is essential when NIP has tag zero.

Every field reader returns a tail no greater than its input. Therefore
one command starting from a positive proof-list value $v$ returns a
proof-list value at most $v-1$, unless it halts. All arithmetic macros
terminate on all natural inputs. The full macro contracts also restore
selector $A=0$ and scratches $X=Y=0$ at every nonhalting command return;
they preserve $N$ outside candidate fetch and leave $L$ equal to the
unconsumed unpairing remainder. At fresh fetch the source loads $N$
as the new candidate and increments $N$ exactly once. Every nonhalting
candidate therefore drains after finitely many commands; the search
then proceeds to the next candidate. In particular, zero selectors or
zero field tails do not cause an omitted case or a nonhalting inner
computation.

Suppose $\Kd\vdash1$, let $\Kblk_1$ be its compiled block, and put
\[
 S=[\Kdrop,\Knip,\Knip],\qquad
 R=\Kenc_\pi(\Kblk_1S).
\]
For any natural padding value $Y$, prepend
$\operatorname{Start}_Y=\Kax(0;1,Y)$. It sets the counter and cache to
zero from every incoming retained tuple. Thus
$[\operatorname{Start}_Y]\Kblk_1S$ halts from every incoming tuple:
the block either halts, or returns top $1$ and zero cache, and $S$
then puts $1$ in the cache and tests it.
Its exact first-candidate number is
\begin{equation}
 N_Y=\Kp\bigl(3,\Kp(Y,\Kp(1,R))\bigr)\geq Y.
 \label{eq:k216-cofinal}
\end{equation}
The inequality follows from the coordinate inequalities for Cantor
pairing. This gives a cofinal family of successful candidates for both
permutations. Choose $Y\geq4$. The successful candidate $N_Y$ is at or
beyond the actual initial next-candidate value $4$. Finite draining of
earlier candidates guarantees that the search reaches it unless a halt
has already occurred. More generally the same proof works beyond every
natural cutoff, and after every finite initial partial candidate.

\paragraph{Source calibration theorem.}
For either source permutation, started at its genuine source root with
register tuple
\[
 (N,T,W,A,D,H,X,Y,L)=(4,0,0,0,0,0,0,0,0),
\]
ZFC proves
\[
 \operatorname{Halt}_{\rm source}
 \quad\Longleftrightarrow\quad\Kd\vdash1
 \quad\Longleftrightarrow\quad
 \exists p\;\operatorname{Proof}_{\mathrm{ZFC}}(p,\bot).
\]
\emph{Proof.} The zero top, zero cache, and all heads of the zero stack
are truth theorems. The soundness invariant proves the forward
implication. The explicit compiler, cofinal family
\eqref{eq:k216-cofinal}, and draining argument prove the reverse
implication. Equation~\eqref{eq:k216-inherited-calibration} gives the
last equivalence. The argument uses finite source identities,
primitive-recursive proof and word transformations, and induction on
natural numbers and finite derivations. It is formalizable in ZFC
without assuming $\operatorname{Con}(\mathrm{ZFC})$.

Combining this theorem with a separate exact binary runtime simulation
and a blank-tape trace to the displayed represented source root yields
\[
 \mathrm{ZFC}\vdash
 \bigl(\text{the corresponding binary machine does not halt on blank tape}\bigr)
 \ \longleftrightarrow\ \operatorname{Con}(\mathrm{ZFC}).
\]
The trace's tape header, head position, gap and spare-word requirements
belong to that binary certificate; they are not consequences of the
source calibration alone.
\endgroup
